% ============================================================================
%  Chapter 22 -- Verm1, Part 4: The Finished Game                (Lecture 22)
% ============================================================================
\chapter{Verm1, Part 4: The Finished Game}
\label{ch:verm1p4}

\epigraphline{Premature optimization is the root of all evil.}{Donald Knuth, 1974}

\begin{chapgoals}
By the end of this chapter you will have
\begin{itemize}
\item the score and the best score written in digits, converted by your own division by ten;
\item a worm that goes faster as the score grows, up to a limit;
\item two sounds: a short note for a meal, a long low one for the end;
\item a \code{draw} function that touches a handful of cells per step instead of the whole screen;
\item the finished \vermone, ready to be packaged in Chapter~\ref{ch:store}.
\end{itemize}
\end{chapgoals}

In this last tutorial on the program the machine itself enters the game.
Knuth's sentence is its motto: the game was made to \emph{work} in Part~3,
and only now, with a measure in hand, is it made \emph{fast}.
Table~\ref{tab:v4:from} lists what is taken from Chapters \ref{ch:processor}
to~\ref{ch:io}; Chapter~\ref{ch:information} comes back too, because a digit
becomes a character by adding 48, the ASCII code of \texttt{0}. The reference
file is \fname{verm1\_part4.S} (Appendix~\ref{app:verm1}).

\begin{table}[htb]
\centering\small
\caption{What Part 4 takes from the chapters on the machine.}\label{tab:v4:from}
\begin{tabularx}{\textwidth}{@{}lLL@{}}
\toprule
\textbf{From} & \textbf{Idea} & \textbf{Where it appears}\\
\midrule
Chapter \ref{ch:processor} & the calling convention: \reg{a}, \reg{t}, \reg{s} registers, \reg{ra}, the stack & small functions with contracts; leaf functions\\
Chapter \ref{ch:processor} & the M extension: \ins{divu} and \ins{remu}, quotient and remainder & the digits of the score\\
Chapter \ref{ch:instructions} & the processor cycle: every instruction costs time & the measure before and after step 4\\
Chapter \ref{ch:video} & the screen is memory: writing pixels is the cost & draw only what has changed\\
Chapter \ref{ch:io} & devices belong to the runtime: the program asks, the runtime does & \code{prg32\_audio\_note}, with a duration\\
\bottomrule
\end{tabularx}
\end{table}

\begin{figure}[htb]
\centering
\begin{subfigure}[b]{0.32\textwidth}\includegraphics[width=\textwidth]{verm1_p4_title.png}\caption{\texttt{TITLE}, with the best score}\end{subfigure}\hfill
\begin{subfigure}[b]{0.32\textwidth}\includegraphics[width=\textwidth]{verm1_p4_play.png}\caption{\texttt{PLAY}: the score in digits}\label{fig:v4:play}\end{subfigure}\hfill
\begin{subfigure}[b]{0.32\textwidth}\includegraphics[width=\textwidth]{verm1_p4_over.png}\caption{\texttt{OVER}}\label{fig:v4:over}\end{subfigure}
\caption{Where this chapter is going.}
\end{figure}

Start from a Part~3 that runs; if yours does not, copy
\fname{verm1\_part3.S}. The chapter has four steps, and the program runs
after each one.

\begin{prgtry}[Before you start: measure]
Run Part~3 in debug mode and write down what the telemetry shows during a
game: retired instructions, frames per second, late frames. The numbers will
be needed in step~5.
\end{prgtry}

% ----------------------------------------------------------------------------
\section{Step 1: the score in digits}

\subsection{A number is not its digits}

The variable \code{score} holds a binary number: 137 is \hex{89}, one byte.
The screen wants three characters, \texttt{'1'}, \texttt{'3'}, \texttt{'7'},
that is, the bytes \hex{31}, \hex{33}, \hex{37}. The conversion is the
algorithm of Chapter~\ref{ch:integers}, done by the program instead of by
hand: divide by ten, three times.
\begin{align*}
137 / 10 &= 13, &&\text{remainder 7: the units;}\\
13 / 10 &= 1, &&\text{remainder 3: the tens;}\\
1 / 10 &= 0, &&\text{remainder 1: the hundreds.}
\end{align*}
The digits come out from the right, and are stored from the last character to
the first. A digit $d$ becomes a character by adding 48.

\subsection{\texttt{to\_decimal}}

\asmfile[caption={Step 1 \textbf{(new)}: \code{to\_decimal}.},label=lst:v4:dec]{code/l22/s08.S}

The contract of Listing~\ref{lst:v4:dec} is: $\reg{a0}=\text{value}$,
$\reg{a1}=\text{where the three characters go}$. \ins{remu} gives the digit
and \ins{divu} what is left: two instructions of the M extension. \ins{sb}
stores a byte at the offsets 2, 1 and 0 from \reg{a1}. The function is a
leaf that uses only \reg{t} registers: it needs no frame at all. Values above
999 would show only their last three digits, a limit that no player will meet
soon. The division instructions are slow, tens of clock cycles each on a
small processor; ten of them per step of the worm are nothing.

\subsection{\texttt{draw\_score}: two numbers, two texts}

\asmfile[caption={Step 1 \textbf{(changed)}: \code{draw\_score}.},label=lst:v4:score]{code/l22/s09.S}

Replace \code{draw\_score} of Part~3 with Listing~\ref{lst:v4:score}: the
marks go. The digits are written into a string in RAM, and then the string is
drawn like any other; the address \code{best\_txt+5} skips the word
\texttt{BEST} and the space, five characters. Text is drawn with its
background colour, so that new digits cover the old ones without any
clearing.

\subsection{The variables of Part 4}

\asmfile[caption={Step 1 \textbf{(changed)}: \code{.data}.},label=lst:v4:data]{code/l22/s10.S}

Replace the \code{.data} of Part~3 with Listing~\ref{lst:v4:data}, and type
all of it now. The two strings \code{score\_txt} and \code{best\_txt} are in
\code{.data} and not in \code{.rodata}, because the program writes into them.
The six variables from \code{dirty} to \code{erase\_y} are not used until
step~3.

\begin{prgcheck}[Run it: digits]
\begin{enumerate}[nosep]
\item Build and run, press \key{A}: the top line reads \texttt{VERM1 000} and,
  on the right, \texttt{BEST 000}.
\item Play and eat: \texttt{001}, \texttt{002}, \ldots
\item Lose: on the game-over screen \texttt{BEST} shows your score.
\item In the debugger find \code{score\_txt} in memory and read its bytes:
  \texttt{30 30 33 00} after three meals.
\end{enumerate}
\end{prgcheck}

% ----------------------------------------------------------------------------
\section{Step 2: levels}

The speed becomes a function of the score.

\asmfile[caption={Step 2 \textbf{(changed)}: constants. Only the two \code{SPEED} lines differ from Part~3.},label=lst:v4:const]{code/l22/s13.S}

\code{SPEED\_Q8} is replaced by a minimum and a maximum
(Listing~\ref{lst:v4:const}): $32/256$, one step every 8 frames, at the
start, and $96/256$, three steps every 8 frames, at most.

\asmfile[caption={Step 2 \textbf{(changed)}: the timer at the end of \code{play\_update}.},label=lst:v4:speed]{code/l22/s14.S}

The speed is $32+4\times\mathit{score}$: a shift and an add
(Listing~\ref{lst:v4:speed}). Then \ins{ble} and \ins{mv} take the smaller of
the speed and 96. It is a \emph{saturation}: the value follows a line and
then stays at the limit, like the clamp of Chapter~\ref{ch:integers}. The
limit is reached at 16 points (Table~\ref{tab:v4:levels}). With fixed point
the speed changes by a sixty-fourth of a cell per frame at each meal: no
integer speed could be so fine.

\begin{table}[htb]
\centering\small
\caption{The levels, in numbers. $256/\mathit{speed}$ gives the frames per
step.}\label{tab:v4:levels}
\begin{tabular}{@{}cccc@{}}
\toprule
\textbf{Score} & \textbf{Speed ($1/256$ cell/frame)} & \textbf{Frames per step} & \textbf{Cells/s at 30 fps}\\
\midrule
0 & 32 & 8 & 3.75\\
4 & 48 & 5.3 & 5.6\\
8 & 64 & 4 & 7.5\\
12 & 80 & 3.2 & 9.4\\
16 and more & 96 & 2.7 & 11.25\\
\bottomrule
\end{tabular}
\end{table}

Fractions of a frame are real: at $48/256$ the steps come every 5 or 6
frames, 5.33 on average. Compute the row of score 2 as an exercise.

\begin{prgcheck}[Run it: faster and faster]
\begin{enumerate}[nosep]
\item Build and run: the worm starts more slowly than in Part~3, and gains
  speed at every meal.
\item Count with the debugger the frames between two steps at score 0: eight.
\item Reach 16 points: the speed no longer changes.
\item Try $\code{SPEED\_MAX}=256$, if you dare.
\end{enumerate}
\end{prgcheck}

% ----------------------------------------------------------------------------
\section{Step 3: \texttt{step} tells what it changed}

This step adds the sounds, and the notes that \code{draw} will need in
step~4.

\subsection{\texttt{play\_note}: one function, any note}

\asmfile[caption={Step 3 \textbf{(changed)}: \code{play\_note} replaces \code{eat\_sound}.},label=lst:v4:note]{code/l22/s18.S}

The caller chooses the note and the duration, in \reg{a2} and \reg{a4}; the
function fills in the other three arguments
(Listing~\ref{lst:v4:note}). The arguments stay in the registers where
\code{prg32\_audio\_note} expects them: designing a contract round the callee
saves moves.

\subsection{\texttt{step}: what it leaves for \texttt{draw}}

The function \code{step} of Part~3 is changed in five places.

\asmfile[caption={Step 3 \textbf{(changed)}: \code{step} remembers the old head.},label=lst:v4:s1]{code/l22/s19.S}

Two stores are added (Listing~\ref{lst:v4:s1}): \code{old\_x} and
\code{old\_y}, the cell that the head is leaving, which \code{draw} will
repaint in the colour of the body.

\asmfile[caption={Step 3 \textbf{(changed)}: \code{step}, when the worm eats.},label=lst:v4:s2]{code/l22/s20.S}

When the worm eats, the tail stays and there is nothing to erase:
$\code{erase\_on}=0$ says so to \code{draw} (Listing~\ref{lst:v4:s2}).

\asmfile[caption={Step 3 \textbf{(changed)}: \code{step} notes the tail, and calls the end.},label=lst:v4:s3]{code/l22/s21.S}

When it does not eat, the tail leaves a cell, and \code{step} remembers
which: $\code{erase\_on}=1$ means that \code{draw} must paint it black
(Listing~\ref{lst:v4:s3}). The end of this fragment is shorter than in
Part~3, because the collision is now one call, to \code{game\_is\_over}.

\asmfile[caption={Step 3 \textbf{(changed)}: \code{step} writes the new head (as in Part~3).},label=lst:v4:s4]{code/l22/s22.S}

\asmfile[caption={Step 3 \textbf{(new)}: the end of \code{step}.},label=lst:v4:s5]{code/l22/s23.S}

At the end (Listing~\ref{lst:v4:s5}) $\code{dirty}=1$ records that a step
has happened and that \code{draw} has work to do. A meal is a point, new food
and a short high note: note 72 for 60 milliseconds, the sound of Part~2
through the new function.

\subsection{\texttt{game\_is\_over}: the transition, in one place}

\asmfile[caption={Step 3 \textbf{(new)}: \code{game\_is\_over}.},label=lst:v4:over]{code/l22/s24.S}

A transition of the state machine deserves a function of its own
(Listing~\ref{lst:v4:over}): the state, the outputs and the sound are in one
place. The function stores \texttt{OVER}, updates the best score as in
Part~3, sets $\code{dirty}=2$ because a new screen must be drawn, and plays a
long low note: 36, for 400 milliseconds.

\begin{prgcheck}[Run it: the sound of defeat]
\begin{enumerate}[nosep]
\item Build and run: the game looks exactly as after step~2.
\item A meal sounds as before; biting yourself gives a long low note.
\item In the debugger watch \code{erase\_x} and \code{erase\_y} follow the
  tail, one step behind.
\item \code{dirty} stays at 1 or 2: nobody reads it yet.
\end{enumerate}
\textbf{Expected result.} No visible change: the screen is still drawn
completely at every frame. The new variables are written and not yet used.
\end{prgcheck}

% ----------------------------------------------------------------------------
\section{Step 4: draw only what changed}

\subsection{What does a frame cost?}

Part~3 clears 64\,000 pixels and redraws the title, the score, the whole worm
and the food, thirty times per second. But between two steps nothing on the
screen changes, and a step changes very little: at most the cell that the
tail has left, which becomes black; the old head, which becomes body; the new
head; and perhaps the food and the score.

\begin{prgconcept}
Remember what changed, and draw only that. The picture stays on the screen by
itself: the frame buffer is memory (Chapter~\ref{ch:video}).
\end{prgconcept}

A worm of 200 cells costs 200 rectangles per frame in Part~3, and three per
step in Part~4. This is what the old home computers did, because they had no
choice.

\subsection{\texttt{draw}: nothing, a little, or everything}

\asmfile[caption={Step 4 \textbf{(changed)}: \code{verm1\_draw}.},label=lst:v4:draw]{code/l22/s28.S}

The variable \code{dirty} has three values (Listing~\ref{lst:v4:draw}). With
0 the function returns at once, with no frame and no call: the common case,
in which nothing has changed, costs five machine instructions. With 1 a step
has happened, and the changes are drawn. With 2 there is a new screen, and
everything is drawn. The flag is cleared as soon as it is read.

\subsection{\texttt{draw\_changes}: three cells and the score}

\asmfile[caption={Step 4 \textbf{(new)}: \code{draw\_changes}.},label=lst:v4:changes]{code/l22/s29.S}

The order matters (Listing~\ref{lst:v4:changes}): first the cell that the
tail has left, if any, in black; then the old head, in the colour of the
body; then the new head, the food and the score. Since the head may enter the
cell that the tail has just left, erasing first and drawing the head
afterwards gives the right picture.

\subsection{\texttt{draw\_full}: a whole screen}

\asmfile[caption={Step 4 \textbf{(new)}: \code{draw\_full}, first half.},label=lst:v4:full1]{code/l22/s30.S}

\asmfile[caption={Step 4 \textbf{(new)}: \code{draw\_full}, \texttt{PLAY} and \texttt{OVER}.},label=lst:v4:full2]{code/l22/s31.S}

\code{draw\_full} is the body of the old \code{verm1\_draw}, and is called
once for each change of state (Listings \ref{lst:v4:full1}
and~\ref{lst:v4:full2}). The score is drawn on the title screen too.

\subsection{Three small functions}

\asmfile[caption={Step 4 \textbf{(new)}: \code{draw\_head}.},label=lst:v4:head]{code/l22/s32.S}

\asmfile[caption={Step 4 \textbf{(new)}: \code{draw\_food}.},label=lst:v4:food]{code/l22/s33.S}

\asmfile[caption={Step 4: \code{draw\_text}, the helper of Part~3.},label=lst:v4:text]{code/l22/s34.S}

\subsection{Every transition asks for a screen}

Three small additions complete the mechanism: wherever the state changes,
\code{dirty} becomes 2.

\asmfile[caption={Step 4 \textbf{(changed)}: \code{verm1\_init} asks for the first screen.},label=lst:v4:init]{code/l22/s35.S}

\asmfile[caption={Step 4 \textbf{(changed)}: \code{verm1\_update}, from \texttt{OVER} to \texttt{TITLE}.},label=lst:v4:upd]{code/l22/s36.S}

\asmfile[caption={Step 4 \textbf{(changed)}: the end of \code{new\_game}, from \texttt{TITLE} to \texttt{PLAY}.},label=lst:v4:new]{code/l22/s37.S}

\begin{prgwarn}
Every transition of the machine asks for a full drawing; every step asks for
a partial one. Forgetting one of them is the typical bug of this technique:
the screen shows a state that is no longer true.
\end{prgwarn}

\begin{prgcheck}[Run it: the finished Verm1]
The same game as after step~3, pixel for pixel (Figures \ref{fig:v4:play}
and~\ref{fig:v4:over}). If you see trails or holes, a change was not
recorded: use the debugger.
\end{prgcheck}

% ----------------------------------------------------------------------------
\section{Step 5: measure}

\prgqt is now used as an instrument: count before you believe.

\begin{prgtry}[Before and after, with numbers]
\begin{enumerate}[nosep]
\item Run the finished game in debug mode and start a game.
\item Read the telemetry during ten seconds of play: retired instructions,
  frames per second, late frames.
\item Compare with the numbers written down for Part~3 at the beginning of the
  chapter.
\item Let the worm grow to twenty cells in both versions and read again: one
  number grows with the worm, the other does not.
\item Put a breakpoint on the first instruction of \code{draw\_changes}: it
  stops once per step, not once per frame.
\item Write the numbers in your notes: they go in the description of the game
  in Chapter~\ref{ch:store}.
\end{enumerate}
\end{prgtry}

The numbers are deliberately not printed here: they must be measured. One
caution in reading them: the telemetry counts the instructions of the
cartridge, not those of the runtime that fills the rectangles, so that the
real saving on the board is larger than the count shows; there every pixel
also travels over SPI to the display.

\subsection{Follow a call with the debugger}

The calling convention of Chapter~\ref{ch:processor} can be watched
instruction by instruction.

\begin{enumerate}
\item Break on \code{call to\_decimal} inside \code{draw\_score}, and read
  \reg{a0}, \reg{a1}, \reg{ra} and \reg{sp}.
\item Step once: the PC jumps into \code{to\_decimal}, and \reg{ra} now holds
  the address of the instruction after the call.
\item Step through: only \reg{t} registers and \reg{a0} change, and \reg{sp}
  never moves. It is a leaf function.
\item At \ins{ret} the PC takes the value of \reg{ra}: the program is back in
  \code{draw\_score}.
\item Do the same with \code{call draw\_text}: \reg{sp} goes down by 16,
  \reg{ra} is stored at \code{12(sp)}, and comes back.
\end{enumerate}

In the listing of the debugger \ins{call} appears as \ins{jal}, or as
\ins{auipc} followed by \ins{jalr}, and \ins{ret} as \code{c.jr ra}.

% ----------------------------------------------------------------------------
\section{Check your understanding}

\begin{quiz}{\code{to\_decimal} receives $\reg{a0}=205$. Which bytes does it store, from offset 0 to 2?}
\item \texttt{02 00 05}
\item \texttt{32 30 35}
\item \texttt{35 30 32}
\item \texttt{cd 00 00}
\end{quiz}

\begin{quiz}{At which score does the worm take one step every 4 frames?}
\item 4
\item 8
\item 16
\item 32
\end{quiz}

\begin{quiz}{\code{dirty} is 0 when \code{verm1\_draw} is called. How many machine instructions does it execute?}
\item 1
\item 5
\item about 50
\item thousands
\end{quiz}

\begin{quiz}{Why can \code{to\_decimal} do without a stack frame?}
\item it is short
\item it calls no function and uses only \reg{a} and \reg{t} registers
\item it has no arguments
\item frames are optional
\end{quiz}

\section{Make it yours}

\begin{exercise}[on \prg]
\begin{enumerate}[nosep]
\item A pause: button \key{B} stops the game and writes \texttt{PAUSE};
  \key{B} again resumes. A fourth state.
\item A melody on the title screen: a table of notes in \code{.rodata}, one
  note every 15 frames.
\item Bonus food: every fifth piece is yellow, is worth three points, and
  disappears after 100 frames.
\item Your own colours with \code{prg32\_palette\_set}, or your own rules:
  this is your game now.
\end{enumerate}
\expected{The version published in Chapter~\ref{ch:store} should have
something that is yours. Change the title string, too. (Hints: the pause is
a state \texttt{PAUSED} with transitions on the edge of \key{B}, \hex{20},
and $\code{dirty}=2$ at both transitions. The melody needs an index and a
frame counter in \code{.data}; the table is read like \code{dx}.)}
\end{exercise}

\section{The functions that did not change}

Eight parts of the file are those of Part~3 and are kept as they are. They
are reprinted here so that this chapter contains the whole program.

\asmfile[caption={\textbf{Unchanged}: \code{verm1\_update}, buttons and edges.},label=lst:v4:u1]{code/l22/s48.S}
\asmfile[caption={\textbf{Unchanged}: \code{new\_game}, the map and the worm.},label=lst:v4:u2]{code/l22/s49.S}
\asmfile[caption={\textbf{Unchanged}: \code{play\_update}, the steering.},label=lst:v4:u3]{code/l22/s50.S}
\asmfile[caption={\textbf{Unchanged}: \code{step}, the table lookup.},label=lst:v4:u4]{code/l22/s51.S}
\asmfile[caption={\textbf{Unchanged}: \code{place\_food}.},label=lst:v4:u5]{code/l22/s52.S}
\asmfile[caption={\textbf{Unchanged}: \code{draw\_worm}, used by \code{draw\_full} only.},label=lst:v4:u6]{code/l22/s53.S}
\asmfile[caption={\textbf{Unchanged}: \code{draw\_cell}.},label=lst:v4:u7]{code/l22/s54.S}
\asmfile[caption={\textbf{Unchanged}: \code{.rodata}.},label=lst:v4:u8]{code/l22/s55.S}

\section{Where we are, and what comes next}

\vermone is finished: about six hundred lines of \riscv assembly, all written
by the reader. It began as three functions that return, and grew one runnable
step at a time. Every chapter of the book is in it: numbers, logic, state,
memory, calls, pixels, devices.

A program that only its author can run is half finished. The next chapter
gives it a name, an icon, credits, a bundle for two architectures, and the
Cartridge Store.
